% Generated draft: requires mathematical and source review.
\documentclass[11pt]{article}
\usepackage[margin=25mm]{geometry}
\usepackage{fontspec}
\setmainfont{Latin Modern Roman}
\usepackage{amsmath,amssymb,mathtools,unicode-math}
\setmathfont{Latin Modern Math}
\usepackage{enumitem,needspace,xurl,hyperref}
\setlength{\parindent}{0pt}
\setlength{\parskip}{5pt}
\setlength{\emergencystretch}{3em}
\newcommand{\N}{\mathbb{N}}
\newcommand{\Z}{\mathbb{Z}}
\newcommand{\Q}{\mathbb{Q}}
\newcommand{\R}{\mathbb{R}}
\newcommand{\C}{\mathbb{C}}
\newcommand{\F}{\mathbb{F}}
\newcommand{\A}{\mathbb{A}}
\newcommand{\PP}{\mathbb{P}}
\newcommand{\E}{\mathbb{E}}
\newcommand{\eps}{\varepsilon}
\newcommand{\dd}{\,\mathrm{d}}
\begin{document}
\section*{UnsolvedMath ID: 30000848}
\textbf{Source problem code:} catalogue entry
\subsection{Definitions and mathematical statement}
Fix a prime \(p\) and a height \(n\geq1\). In the usual Morava \(E\)-theory convention let \(R=W(\F_{p^n})\), the Witt ring of the finite field, and
\[
 E_{n*}\cong R[[u_1,\ldots,u_{n-1}]][u^{\pm1}],
 \qquad |u_i|=0,\quad |u|=-2.
\]
Let \(\mathbb S_n\) be the Morava stabilizer group of the height-\(n\) formal group, with its action on these coefficients, and let \(G\leq\mathbb S_n\) be finite. The source writes the coefficient ring as \(\mathbb Z_{p^n}\); interpreting this as \(W(\F_{p^n})\), rather than \(\Z/p^n\Z\), is the standard unramified-coefficient interpretation and requires confirmation of the report's notation.

Write \(M_n\) for the Dieudonn\'e module named in the original statement, equipped with the intended \(G\)-action and placed in degree \(-2\), and form the graded commutative symmetric algebra \(S_R(M_n)=\operatorname{Sym}_R(M_n)\). For a specified vector \(m\in M_n\), a multiplicative norm class would have the form
\[
 N_G(m)=\prod_{g\in G}g\cdot m\ \in S_R(M_n)^G.
\]
The conjecture uses the \(G\)-invariant norm class \(\Delta\) intended in the report; the extraction does not specify a vector or another construction selecting that class. Set \(A=S_R(M_n)[\Delta^{-1}]\). For the intended degree-zero ideal \(I\subseteq A_0\), completion means \(\widehat A_I=\varprojlim_k A/I^kA\), with the grading/degreewise convention to be fixed from the source. The ideal must be compatible with the \(G\)-action to obtain the indicated action on the completion.

The conjecture attributed to Hopkins asserts, for every such finite subgroup \(G\), a \(G\)-equivariant isomorphism of the intended graded coefficient algebras
\[
 E_{n*}\ \cong\ \bigl(S_{\mathbb Z_{p^n}}(M_n)[\Delta^{-1}]\bigr)^{\wedge}_I.
\]
This is a presentation of coefficient algebras, not an asserted equivalence of spectra. The source does not mean completion at every arbitrary degree-zero ideal: the relevant \(I\), module, norm class, and equivariant algebra structures must be those in the report.

The original uses \(E_n^*\), whereas the cleaned statement uses \(E_{n*}\) and specifies degree \(-2\); homological versus cohomological grading must be reconciled. The supplied assessment dated 2026-08-21 says the extraction omits an unnamed conjecture's statement, but the supplied original, clean statement, and verification note explicitly state and attribute the Hopkins conjecture. That assessment conflicts with the present record and is not evidence about this conjecture's resolution.
\subsection{Short English statement}
Hopkins' conjecture seeks an equivariant symmetric-algebra presentation of Morava \(E\)-theory coefficients after inverting a norm class and completing. The module, norm element, ideal, coefficient-ring notation, and grading need source-level specification. The dated assessment appears inconsistent with the actual statement supplied here.
\subsection{Sources}
\begin{itemize}
\item Local UnsolvedMath registry, permanent ID \(30000848\), exact source code \texttt{catalogue entry}, proposed year 2007. The original explicitly attributes the conjecture to Hopkins; no separate proposer field is supplied.
\item \emph{Homotopy Theory (2008)}, Oberwolfach Reports, DOI 10.4171/owr/2007/45, pp. 2671--2726. \url{https://doi.org/10.4171/owr/2007/45}.
\item The supplied \texttt{corrected\_verified} note says the official report states the Hopkins conjecture and defines its symmetric algebra, grading, norm class, localization, and completion. Those full definitions are not reproduced in the extracted record.
\item The open label and literature assessment checked 2026-08-21 are supplied dated triage. Its claim that only motivation for an unnamed conjecture is present conflicts with the explicit original and cleaned statements; no independent status or source assessment is claimed here.
\end{itemize}
\end{document}
